\title{Direct Preference Optimization: Your Language Model is Secretly a Reward Model}

\begin{abstract}
Although large unsupervised language models (LMs) acquire extensive world knowledge and demonstrate some capability for reasoning, exerting precise control over their output remains challenging because of their entirely unsupervised training paradigms. Current approaches to enhance model controllability typically involve gathering human annotations comparing the relative merits of model outputs, and subsequently adapting the LM to mirror these human preferences—most frequently via reinforcement learning from human feedback (RLHF). Yet, RLHF introduces significant complexity and instability by first fitting a reward model to represent human judgment, followed by leveraging reinforcement learning to optimize the LM for this reward, all while striving to prevent deviation from the pre-trained model’s behavior. In this work, we propose a novel way of parameterizing the reward model within the RLHF setup that facilitates direct derivation of the associated optimal policy in closed form. This leads to a reformulation of the standard RLHF problem as a tractable classification task. We introduce \textit{Direct Preference Optimization} (DPO), an approach that is robust, effective, and computationally efficient, obviating the need for either LM sampling during fine-tuning or extensive hyperparameter searches. Empirical results demonstrate that DPO can adapt LMs to follow human preferences as well as, or even surpassing, existing techniques. Remarkably, DPO outperforms PPO-based RLHF in modulating output sentiment and achieves equivalent or superior results for summarization and single-turn dialogue tasks, while significantly reducing implementation and training complexity.
\end{abstract}

\section{Introduction}
Large-scale unsupervised language models (LMs), when exposed to vast corpora of data, develop remarkable and sometimes unexpected capabilities~\citep{chowdhery2022palm, brown2020language, touvron2023llama,bubeck2023sparks}. Nevertheless, these models are trained on human-generated data that reflects a broad spectrum of intentions, priorities, and expertise. Not all such human attributes are beneficial to replicate; for instance, although it is advantageous for an AI coding assistant to \textit{recognize} typical programming errors to correct them, ideally, the model’s outputs should favor the infrequent but high-quality programming skills embedded in the data. Similarly, while awareness of a misconception held by half the population may be useful, it is undesirable for the model to endorse this misconception as true in half of relevant responses. Therefore, determining the model's \emph{intended outputs and conduct} from among its expansive \textit{knowledge base and competencies} is essential for building AI systems that are reliable, effective, and amenable to control \citep{ouyang2022training}. Although reinforcement learning (RL) is commonly employed to align LMs with human values, we demonstrate that the RL-derived objective prevalent in these methods can in fact be exactly solved via a straightforward binary cross-entropy loss, offering significant simplification to the preference optimization process.

Broadly, current approaches impart preferred behaviors into language models by leveraging human preference datasets designed to capture the safety and helpfulness standards desired by users. This preference alignment is conducted after the preliminary, unsupervised phase, during which the LM is pretrained on large-scale text. While one could use supervised fine-tuning on annotated demonstrations that exemplify desirable responses, reinforcement learning from human (or AI) feedback (RLHF/RLAIF; \citep{christiano2017deep,bai2022constitutional}) represents the dominant methodology. RLHF consists of training a reward model on human preference data and then applying RL to adjust the model’s policy so that its outputs are scored favorably, all while maintaining proximity to the original parameters. Although RLHF has enabled LMs to achieve outstanding performance in both conversation and code generation, this approach is much more intricate than supervised techniques, as it demands the training of several language models and policy sampling throughout the training cycle, resulting in considerable computational expense.

In this work, we demonstrate a method for optimizing language models in alignment with human preferences without relying on explicit reward modeling or reinforcement learning frameworks. We introduce \textit{{Direct Preference Optimization} (DPO)}, a procedure that tacitly targets the same objective as standard RLHF approaches—namely, maximizing rewards subject to a KL-divergence constraint—while offering a more straightforward and easily implemented solution. Conceptually, the DPO mechanism adjusts model parameters by enhancing the log odds in favor of human-preferred responses over those less favored. Critically, it employs dynamically assigned, instance-level importance weights, thereby mitigating the model collapse observed when utilizing a naïve likelihood ratio formulation. DPO, akin to prior algorithms, draws upon theoretical constructs such as the Bradley-Terry model (\cite{bradley1952rankanalysis}) to gauge how closely a reward function tracks with human preference annotations. The distinction lies in DPO’s use of a variable transformation: the preference loss is articulated directly as a function of the policy rather than through an intermediate reward model. This enables the training of a language model policy using a simple binary cross entropy objective directly over preference-labeled examples, resulting in the optimal policy for a latent reward signal derived from preference data.

The central advancement presented in this paper is Direct Preference Optimization (DPO): an RL-independent, streamlined method for preference-driven training of language models. Our empirical results indicate that DPO matches or surpasses the efficacy of prevalent techniques such as PPO-based RLHF in learning from preference data on tasks like sentiment adjustment, abstractive summarization, and conversational modeling, evaluated on models with up to 6B parameters.

\section{Related Work}

Language models trained in a self-supervised manner and scaled up have demonstrated the ability to handle certain tasks in a zero-shot setting \citep{radford2019language}, or with few-shot examples provided in prompts \citep{gpt3,megatron,chowdhery2022palm}. Nevertheless, the effectiveness of these models on downstream applications and their alignment with user intentions can be greatly enhanced by fine-tuning them on curated collections of instructions and human-generated completions \citep{mishra-etal-2022-cross,sanh2022multitask,chung2022scaling,thoppilan2022lamda}. This process, referred to as `instruction-tuning,’ enables large language models (LLMs) to extrapolate to new instructions beyond those seen during fine-tuning, thereby improving their overall utility \citep{chung2022scaling}. Although instruction-tuning has proven highly effective, gathering \textit{relative} assessments of output quality from humans is often more straightforward than collecting expert-authored demonstrations. As a result, recent research has focused on further adapting LLMs with data reflecting human preferences, leading to gains in areas like translation \citep{kreutzer-etal-2018-reliability}, summarization \citep{stiennon2022learning,ziegler2020finetuning}, narrative generation \citep{ziegler2020finetuning}, and following instructions \citep{ouyang2022training,ramamurthy2023is}. Typically, these approaches first train a neural reward model to be consistent with human preference data, using frameworks such as the Bradley-Terry model \citep{bradley1952rankanalysis}. The language model is then further fine-tuned to maximize this learned reward through reinforcement learning techniques, most notably REINFORCE \citep{williams1992reinforce}, proximal policy optimization (PPO; \cite{schulman2017proximal}), or related algorithms \citep{ramamurthy2023is}. A parallel stream of research utilizes LLMs, already adapted via instruction-tuning with human feedback, to produce synthetic datasets for preferences over specific attributes (such as safety or harmlessness) \citep{bai2022constitutional}. In this context, humans provide only light supervision, usually in the form of a rubric that guides the LLM in annotating data. Collectively, these strategies blend advances in reinforcement learning-based training for language models \citep{Ranzato2015SequenceLT,paulus2018a,wu2018learning} with broad approaches for learning from human feedback \citep{christiano2017deep,kupcsik2018learning}. Despite the attractiveness of leveraging human preference comparisons, reinforcement learning-based fine-tuning of LLMs remains difficult in practice; in this work, we present a principled method for optimizing over relative preferences that bypasses the need for reinforcement learning.

Learning from preferences, outside of language tasks, has been explored in both bandit and reinforcement learning contexts, leading to the development of a variety of methods. In the case of contextual bandit learning, where actions are selected based on preferences or rankings instead of explicit rewards, the problem is referred to as a contextual dueling bandit (CDB; \cite{yue2012karmed,dudik2015contextual}). Since absolute reward information is missing in CDB scenarios, theoretical treatment replaces the standard concept of an optimal policy with the notion of a \textit{von Neumann winner}—defined as a policy that defeats or ties with any alternative policy at least half the time in expectation \citep{dudik2015contextual}. A key distinction in CDBs is that preference feedback is provided in an online manner, contrasting with the offline, fixed batch of human preference annotations typically available in human preference learning \citep{yan2022human}. Likewise, in \textit{preference-based RL} (PbRL), agents receive binary preference signals from an unknown latent `scoring' mechanism instead of reward signals \citep{BusaFekete2014,ruiz2023dueling}. PbRL algorithms, which may leverage off-policy preference data, commonly involve the explicit estimation of this underlying scoring function (the reward model) before proceeding to optimize it \citep{jain2013learning,BusaFekete2014,christiano2017deep,sadigh2017active,kupcsik2018learning}. Contrasting with these multi-stage approaches, we propose an integrated policy learning method that directly optimizes a policy to fulfill observed preferences.

\section{Preliminaries}\label{section:prelims}

We briefly review the RLHF methodology as described by \citeauthor{ziegler2020finetuning} and expanded in \citep{stiennon2022learning, bai2022training, ouyang2022training}. The typical RLHF workflow consists of three main steps: (1) supervised fine-tuning (SFT); (2) collecting preferences and training a reward model; and (3) reinforcement learning for policy optimization.

\textbf{SFT}: The RLHF pipeline typically initiates by further training a pre-existing language model via supervised learning on curated, high-quality datasets tailored to downstream tasks (such as dialogue, summarization, etc.), resulting in a model denoted as $\pi^\text{SFT}$.

\textbf{Reward Modelling Phase}: Next, the fine-tuned SFT model generates outputs by responding to input prompts $x$, producing answer pairs $(y_1, y_2)\sim \pi^\text{SFT}(y \mid x)$. Human annotators then evaluate these pairs, expressing a preference for one response over the other, indicated by $y_w\succ y_l \mid x$, where $y_w$ stands for the preferred completion and $y_l$ the less favored, drawn from $(y_1, y_2)$. These observed preferences are assumed to stem from an underlying reward model $r^*(y, x)$ that is inaccessible during training. A variety of preference modeling approaches exist; a widely used option is the Bradley-Terry (BT) \cite{bradley1952rankanalysis} model, although other ranking models such as Plackett-Luce \citep{plackett1975analysis, luce2012individual} are suitable when preference annotations include full rankings. Under the BT model, the human preference distribution $p^*$ is defined as:

Given access to a fixed set of preference comparisons $\mathcal{D}=\bigl\{x^{(i)}, y_w^{(i)}, y_l^{(i)}\bigr\}_{i=1}^N$ sampled according to $p^*$, one can parameterize a reward model $r_{\phi}(x, y)$ and fit its parameters by maximizing likelihood. If we treat this as a binary classification task, the associated loss function is the negative log-likelihood:

\begin{equation}\label{eq:reward_model}
    \mathcal{L}_R(r_{\phi}, \mathcal{D}) = -\mathbb{E}_{(x, y_w, y_l)\sim \mathcal{D}}\bigl[\log \sigma(r_{\phi}(x, y_w)- r_{\phi}(x, y_l))\bigr]
\end{equation}

where $\sigma$ denotes the logistic (sigmoid) function. For language model applications, $r_{\phi}(x, y)$ is typically initialized from the SFT model $\pi^\text{SFT}(y \mid x)$, and a linear output head is added atop the final transformer layer to yield a scalar reward prediction, as outlined in \cite{ziegler2020finetuning}. To minimize variance in the reward estimates, previous research has normalized the reward outputs so that $\mathbb{E}_{x,y\sim \mathcal{D}}\left[r_\phi(x, y)\right] = 0$ over the dataset.

\textbf{RL Fine-Tuning Phase}: During the reinforcement learning stage, the trained reward function provides feedback for language model optimization. Drawing on methods established in prior work~\citep{jaques2017sequence, jaques2020human}, this procedure seeks to solve:

\begin{equation}\label{eq:RL}
\max_{\pi_{\theta}}  \mathbb{E}_{x\sim \mathcal{D}, y\sim \pi_{\theta}(y \mid x)}\bigl[r_{\phi}(x, y)\bigr] - \beta\mathbb{D}_{\textrm{KL}}\bigl[\pi_{\theta}(y\mid x)\mid \mid \pi_\text{ref}(y\mid x)\bigr],
\end{equation}

where $\beta$ modulates how much the learned policy may diverge from a reference policy $\pi_\text{ref}$, often taken as the original SFT policy $\pi^\text{SFT}$. Common practice initializes the policy $\pi_\theta$ with $\pi^\text{SFT}$ as well. This KL regularization is crucial—it ensures that the optimized policy remains close to the distribution where the reward model is reliable, and also prevents loss of diversity or convergence to trivial deterministic outputs. Since text generation involves discrete actions, this objective is non-differentiable and typically requires reinforcement learning for optimization. Established techniques \citep{ziegler2020finetuning, stiennon2022learning, bai2022training, ouyang2022training} formulate the reward as ${r(x, y) = r_{\phi}(x, y) -\beta (\log \pi_{\theta}(y\mid x) - \log \pi_\text{ref}(y\mid x))}$ and employ the PPO algorithm~\cite{schulman2017proximal} for policy optimization.

\section{Direct Preference Optimization}\label{sec:DPO}

Confronted with the difficulties of scaling RL algorithms for language model fine-tuning, our objective is to introduce a streamlined method for preference-based policy optimization. In contrast to conventional RLHF pipelines—which first infer a reward and subsequently maximize it using RL—our method adopts a distinct reward parameterization that allows direct computation of the optimal policy, thus circumventing the RL training loop entirely. Central to our approach is the exploitation of a precise mapping from reward models to their optimal policies. This analytic correspondence allows us to convert objectives defined over reward functions into analogous objectives over policy parameters. By employing this transformation, we eliminate the need for an explicitly trained standalone reward model, yet still maintain compatibility with standard human preference models such as Bradley-Terry. Ultimately, this framework unifies the representation of the policy and the (implicit) reward within a single network.

\textbf{DPO Objective Derivation.} We begin by considering the reinforcement learning objective found in previous studies, as presented in Eq.~\ref{eq:RL}, for an arbitrary reward function $r$. Consistent with earlier work~\citep{peters2007reinforcement, peng2019advantage, korbak2022reinforcement, go2023aligning}, one can directly verify that the solution that optimizes the KL-regularized reward maximization problem in Eq.~\ref{eq:RL} is given by:

\begin{equation}\label{eq:op_policy}
    \pi_r(y\mid x) = \frac{1}{Z(x)}\pi_\text{ref}(y\mid x)\exp\left(\frac{1}{\beta}r(x, y)\right),
\end{equation}

where $Z(x) =\sum_{y}\pi_\text{ref}(y\mid x)\exp\left(\frac{1}{\beta}r(x, y)\right)$ acts as the normalization constant. The derivation details are included in Appendix \ref{app:derivation1}. Importantly, even when employing a maximum likelihood estimate $r_{\phi}$ for the true reward $r^*$, calculating $Z(x)$ remains computationally intensive~\citep{korbak2022reinforcement, go2023aligning}, thereby making direct utilization of this parameterization impractical. Nonetheless, Eq.~\ref{eq:op_policy} can be rearranged so that the reward function is expressed as a function of the optimal policy $\pi_r$, reference policy $\pi_\text{ref}$, and $Z(\cdot)$. By taking the logarithm of both sides of Eq.~\ref{eq:op_policy} and rearranging terms, we arrive at:

\begin{equation}\label{eq:main_eq}
    r(x,y) =\beta \log \frac{\pi_r(y\mid x)}{\pi_\text{ref}(y\mid x)} + \beta \log Z(x).
\end{equation}

Applying this change of variables to the ground-truth reward $r^*$ and the associated optimal policy $\pi^*$, we observe that the Bradley-Terry model relies solely on reward differences between completions: specifically, ${p^*(y_1 \succ y_2 \mid x) = \sigma(r^*(x, y_1) - r^*(x, y_2))}$. Plugging the parameterization from Eq.~\ref{eq:main_eq} into the preference model Eq.~\ref{eq:bradley-terry}, the partition functions $Z(x)$ eliminate due to subtraction, and human preference probabilities can now be articulated exclusively in terms of the optimal policy $\pi^*$ and reference policy $\pi_\text{ref}$. Therefore, under the Bradley-Terry assumption, the optimal RLHF policy $\pi^*$ fulfills:

\begin{equation}\label{eq:objective}
    p^*(y_1\succ y_2 \mid x)=\frac{1}{1 + \exp\left(\beta \log \frac{\pi^*(y_2\mid x)}{\pi_\text{ref}(y_2\mid x)} - \beta \log \frac{\pi^*(y_1\mid x)}{\pi_\text{ref}(y_1\mid x)}\right)}
\end{equation}

For a comprehensive proof, see Appendix~\ref{app:derivation2}. Although this derivation utilizes the Bradley-Terry formulation, an analogous process yields results for general Plackett-Luce models~\citep{plackett1975analysis, luce2012individual}, discussed in Appendix~\ref{app:plackett_luce_models}.

Having thus represented human preference probabilities through the optimal policy instead of the reward function, we can now develop a maximum likelihood training objective for a parameterized policy $\pi_\theta$. In analogy with the standard reward modeling method (cf. Eq.~\ref{eq:reward_model}), the policy learning objective is given by:

\begin{equation}\label{eq:optimum_model}
    \mathcal{L}_\text{DPO}(\pi_{\theta}; \pi_\text{ref}) = -\mathbb{E}_{(x, y_w, y_l)\sim \mathcal{D}}\left[\log \sigma \left(\beta \log \frac{\pi_{\theta}(y_w\mid x)}{\pi_\text{ref}(y_w\mid x)} - \beta \log \frac{\pi_{\theta}(y_l\mid x)}{\pi_\text{ref}(y_l

In this formulation, we infer an implicit reward by adopting an alternative parameterization, where the optimal policy coincides with $\pi_\theta$. Additionally, since this approach corresponds to optimizing a reparametrized Bradley-Terry model, it possesses certain desirable theoretical guarantees, such as consistency under appropriate assumptions on the distribution of preference data \cite{bong2022generalized}. Section~\ref{sec:theory} provides a more detailed examination of the theoretical aspects of DPO and contrasts them with related literature.

\textbf{Mechanics of the DPO update.} To gain deeper insight into the operation of DPO, it is instructive to study the gradient of its objective, $\mathcal{L}_\text{DPO}$. The gradient with respect to the parameter vector $\theta$ is given by:

\begin{multline*}\label{eq:gradient}
    \nabla_\theta \mathcal{L}_\text{DPO}(\pi_\theta;\pi_\text{ref}) = \\ -\beta\mathbb{E}_{(x, y_w, y_l) \sim \mathcal{D}} \bigg[\underbrace{\sigma(\hat{r}_\theta(x, y_l) - \hat{r}_\theta (x, y_w))}_\text{increased weight when reward order is mistaken}\bigg[\underbrace{\nabla_\theta\log \pi(y_w \mid x)}_\text{promote $y_w$} - \underbrace{\nabla_\theta\log\pi(y_l \mid x)}_\text{penalize $y_l$}\bigg]\bigg],
\end{multline*}

where the implicit reward assigned by the model is $\hat{r}_\theta(x, y) = \beta \log \frac{\pi_\theta(y \mid x)}{\pi_\text{ref}(y \mid x)}$, using the outputs of both the language model $\pi_\theta$ and the reference model $\pi_\text{ref}$ (see Section~\ref{sec:theory} for further elaboration). Conceptually, the gradient of $\mathcal{L}_\text{DPO}$ encourages increasing the likelihood of completions labeled as preferred, $y_w$, while simultaneously reducing the probability assigned to non-preferred responses, $y_l$. The relative influence of each example is modulated by the degree to which the implicit reward model $\hat{r}_\theta$ erroneously prefers $y_l$ over $y_w$, scaled by $\beta$, effectively measuring how much the model misorders the completions and factoring in the magnitude of the KL constraint. Empirically, we observe that this particular weighting is crucial: omitting it (as in a simplistic variant) can lead to detrimental model behavior, such as language model degeneration (see Appendix Table~\ref{tab:unlikelihood_generations} for evidence).

\textbf{DPO overview.} The standard workflow for DPO consists of: 1) Drawing completions $y_1, y_2 \sim \pi_\text{ref}(\cdot \mid x)$ for each prompt $x$, annotating these samples with human preferences to build an offline preference dataset $\mathcal{D} = \{x^{(i)}, y_w^{(i)}, y_l^{(i)}\}_{i=1}^N$; and 2) updating the language model $\pi_\theta$ by minimizing $\mathcal{L}_\text{DPO}$, using the chosen $\pi_\text{ref}$, the dataset $\mathcal{D}$, and target value of $\beta$. 
In practice, rather than generating new completions and collecting fresh human preferences, one typically prefers to leverage existing public preference datasets. As such, since these datasets are usually collected from samples drawn via $\pi^\text{SFT}$, we set $\pi_\text{ref} = \pi^\text{SFT}$ whenever it is accessible. In the absence of $\pi^\text{SFT}$, we instead construct $\pi_\text{ref}$ by maximizing the likelihood over preferred responses ${(x, y_w)}$, i.e., by setting ${\pi_\text{ref} = \argmax_{\pi}\mathbb{E}_{x, y_w \sim \mathcal{D}}\left[\log \pi(y_w \mid x)\right]}$. This initialization helps to address the mismatch between the unknown ideal reference distribution and the practical $\pi_\text{ref}$ employed within DPO. Supplementary information regarding implementation specifics and hyperparameter configurations can be found in Appendix~\ref{app:implementation}.

\section{Theoretical Analysis of DPO} This section aims to deepen the understanding of the DPO approach, offering theoretical justifications, and comparing the benefits of DPO against the shortcomings of actor-critic methods used in RLHF—such as PPO~\cite{schulman2017proximal}.

\label{sec:theory}

\subsection{Your Language Model Is Secretly a Reward Model} DPO manages to avoid both the separate reward modeling step and the need for reinforcement learning to train the policy by using a unified maximum likelihood framework. The loss function in Eq.~\ref{eq:main_eq} aligns with a Bradley-Terry model parameterized by rewards $r^*(x, y) = \beta \log\frac{\pi^*_\theta(y \mid x)}{\pi_\text{ref}(y \mid x)}$, and the training of our model $\pi_{\theta}$ corresponds to reward model learning as in Eq.~\ref{eq:reward_model}, subject to a reparameterization. In what follows, we elaborate the theoretical rationale for this formulation, demonstrate that it retains the full representational power of reward models, and show it enables precise recovery of the optimal policy. Our analysis begins by introducing an equivalence relation for reward functions.

\begin{definition}
Two reward functions $r(x, y)$ and $r'(x, y)$ are defined as equivalent if and only if there exists a function $f$ such that ${r(x, y)-r'(x, y) = f(x)}$.     
\end{definition}

This equivalence can readily be verified as an equivalence relation, which naturally partitions the space of reward functions into equivalence classes. The following two lemmas capture important properties of these classes:

\begin{lemma}\label{lemma:same_prefrence} Within the Plackett-Luce framework, and its special case, the Bradley-Terry model, any pair of reward functions belonging to the same equivalence class will generate identical preference distributions.
\end{lemma}

\begin{lemma}\label{lemma:same_policy}
    Any two reward functions from an equivalence class will yield the same optimal policy when applied to the constrained RL problem.
\end{lemma}

We provide simple proofs of these lemmas in Appendix \ref{app:lemma1}. The first lemma highlights a well-known identifiability issue that arises in the Plackett-Luce family \cite{plackett1975analysis}: different reward functions related by the equivalence definition are observationally indistinguishable. Therefore, additional constraints are often necessary to ensure identifiability of the MLE solutions in Eq. \ref{eq:reward_model} \cite{bong2022generalized}. The second lemma implies that, for our ultimate purposes, it is sufficient to recover any reward function within the optimal equivalence class, as all yield the same optimal policy. In Appendix~\ref{app:thm1}, we formalize this insight in the following theorem:

\begin{theorem}\label{thm:main}
    Assuming mild conditions, all equivalence classes of rewards compatible with Plackett-Luce (and, in particular, Bradley-Terry) models can be expressed via the reparameterization ${r(x, y) = \beta \log \frac{\pi(y\mid x)}{\pi_\text{ref}(y\mid x)}}$, for some model $\pi(y\mid x)$ and a chosen reference $\pi_\text{ref}(y \mid x)$.
\end{theorem}

\begin{sproof}
    Let $r(x, y)$ be an arbitrary reward function with a corresponding optimal policy $\pi_r(y \mid x)$ given by Eq. \ref{eq:op_policy}. To show that a member of $r$'s equivalence class admits the reparameterization in question, define the transformation
\begin{equation}
    f(r; \pi_\text{ref}, \beta)(x, y) = r(x, y) - \beta\log\sum_{y}\pi_\text{ref}(y\mid x)\exp\left(\frac{1}{\beta}r(x, y)\right)
\end{equation}
This mapping adjusts the reward by subtracting the log-partition function (with respect to $\pi_\text{ref}$) and thus only introduces an $x$-dependent offset, keeping it within the same equivalence class as $r(x, y)$. Substituting $r$ with the expression on the right-hand side of Eq.~\ref{eq:main_eq}, valid for any reward, we obtain $f(r; \pi_\text{ref}, \beta)(x, y) = \beta \log \frac{\pi_r(y\mid x)}{\pi_\text{ref}(y\mid x)}$. Thus, this construction yields a reward of the required form from the equivalence class of $r$, ensuring that our reparameterized model is fully expressive up to equivalence.
\end{sproof}

An alternative interpretation of Theorem~\ref{thm:main} is that it identifies the precise reward function chosen by the DPO reparameterization from within each equivalence class, namely, the reward function that fulfills the condition:

\begin{equation}\label{eq:lag_p}
     \sum_{y}\underbrace{\pi_\text{ref}(y\mid x)\exp\left(\frac{1}{\beta}r(x, y)\right)}_{=\pi(y\mid x)\text{, using Thm.~\ref{thm:main} reparam.}} = 1,
\end{equation}

In other words, this guarantees that $\pi(y\mid x)$ forms a legitimate probability distribution, being both positive and normalized to sum to one. Observing Eq.~\ref{eq:op_policy} reveals that Eq.~\ref{eq:lag_p} constitutes the partition function corresponding to the optimal policy generated by the reward function $r(x, y)$. The central idea underlying the DPO algorithm is that it introduces specific constraints on the otherwise underdetermined Plackett-Luce family of preference models (including, in particular, the Bradley-Terry model). This preserves the set of reward models that can be represented, but crucially ensures that the optimal policy from Eq.~\ref{eq:op_policy} remains analytically manageable for every prompt $x$.

\subsection{Instability of Actor-Critic Algorithms} Our framework can further serve to illuminate the sources of instability frequently observed with traditional actor-critic approaches employed in RLHF pipelines, such as PPO. Focusing on the reinforcement learning fine-tuning phase detailed in Section \ref{section:prelims}, connections can be established with the control as inference paradigm \cite{levine2018reinforcement} as applied to the constrained RL formulation described in \ref{eq:RL}. Here, we posit a parameterized policy $\pi_{\theta}(y\mid x)$ and aim to minimize $\mathbb{D}_{\text{KL}}[\pi_{\theta}(y|x) \mid \mid \pi^*(y\mid x)]$, where the reference policy $\pi^*$—defined in Eq.~\ref{eq:optimum_model}—is induced by the reward function $r_{\phi}(y, x)$. Elementary manipulation yields the following maximization objective:

\begin{equation}\label{eq:AC}
    \max_{\pi_{\theta}}\mathbb{E}_{\pi_{\theta}(y\mid x)}\bigg[\underbrace{r_{\phi}(x, y) -\beta\log\sum_{y}\pi_\text{ref}(y\mid x)\exp\left(\frac{1}{\beta}r_{\phi}(x, y)\right)}_{f(r_{\phi}, \pi_\text{ref}, \beta)} - \underbrace{\beta\log\frac{\pi_{\theta}(y\mid x)}{\pi_\text{ref}(y\mid x)}}_{\text{KL}}\bigg]
\end{equation}

The objective considered here is identical to that optimized in previous studies 
\citep{ziegler2020finetuning, stiennon2022learning, bai2022training, ouyang2022training}, which utilize the DPO-equivalent reward for the reward class $r_{\phi}$. Within this framework, the normalization component in $f(r_{\phi}, \pi_\text{ref}, \beta)$ can be interpreted as the soft value function for the reference policy $\pi_\text{ref}$. Although this normalization does not change the optimal policy, omitting it may result in a high-variance policy gradient, thereby reducing training stability. One can introduce a learned value function to approximate the normalization term; however, optimizing such a function can be challenging. Alternatively, earlier approaches often normalize rewards with respect to a human completion baseline, effectively producing a one-sample Monte Carlo estimate of the normalization factor. By contrast, the DPO reparameterization produces a reward function that eliminates the need for any baseline normalization.

\section{Experiments}
Here, we empirically assess the effectiveness of DPO for learning policies solely from preference feedback. Initially, in a controlled text generation environment, we investigate DPO’s efficiency in balancing reward maximization and KL-divergence minimization with the reference policy, and benchmark its performance against established preference-based algorithms like PPO. Subsequently, we extend our evaluation to encompass larger models and more complex RLHF scenarios, such as summarization and dialogue tasks. Our results indicate that, even with minimal hyperparameter adjustment, DPO often matches or surpasses strong baselines like RLHF with PPO and outperforms methods that select the best out of $N$ sampled rollouts under a learned reward model. The experimental design is detailed in the following section, with further information available in Appendix~\ref{app:exp_details}.

\textbf{Tasks.} Our study investigates three open-ended text generation tasks. For every experiment, the methods are trained to derive a policy using a preference dataset $\mathcal{D}=\bigl\{x^{(i)}, y_w^{(i)}, y_l^{(i)}\bigr\}_{i=1}^N$. The first task, \textbf{controlled sentiment generation}, uses $x$ as an initial segment from a movie review within the IMDb dataset \cite{maas-EtAl:2011:ACL-HLT2011}, requiring the policy to produce $y$ that expresses positive sentiment. To enable objective comparison, preference pairs between outputs are automatically produced using a pre-trained sentiment classifier, enforcing $p(\text{positive}\mid x,y_w)>p(\text{positive}\mid x,y_l)$. For SFT, we continue training GPT-2-large on IMDB train set reviews until convergence (details in App~\ref{app:sentiment_details}). In the \textbf{summarization} task, $x$ is a Reddit forum post, and the model must generate a summary $y$ that captures the main points. As in previous work, we adopt the Reddit TL;DR summarization corpus \citep{volske-etal-2017-tl} as well as human preference data collected by \citeauthor{stiennon2022learning}. Our SFT baseline is trained on human-generated forum post summaries\footnote{\url{https://huggingface.co/CarperAI/openai_summarize_tldr_sft}} and implemented with the TRLX \citep{leandro_von_werra_2023_7790115} library for RLHF. Human preferences in this dataset were assembled by \citeauthor{stiennon2022learning} from samples by another SFT model with similar training procedures. The last task, \textbf{single-turn dialogue}, defines $x$ as a user prompt, which can range from scientific queries to personal advice requests. Here, the policy must output $y$ that is informative and user-oriented, using the Anthropic Helpful and Harmless dialogue collection \citep{bai2022training}, which includes 170,000 exchanges between people and an automated system. Each dialogue transcript concludes with two model-generated responses and a human annotation indicating preference. Since a suitable SFT model is unavailable for this scenario, we instead train an out-of-the-box language model only on human-preferred completions to establish the SFT baseline.

\textbf{Evaluation.} We employ two distinct strategies for evaluation in our experiments. To assess each algorithm's proficiency in maximizing constrained reward objectives, we analyze their performance by tracing the trade-off frontier between obtained reward and KL-divergence from the reference policy in the controlled sentiment generation scenario. This is feasible since the ground-truth reward function (a sentiment classifier) is accessible. In contrast, because real-world applications lack such an oracle reward function, we assess algorithms using their \textit{win rate} relative to a baseline policy, relying on GPT-4 as a stand-in for human judgment. For summarization, GPT-4 evaluates summary quality against reference summaries, while for single-turn dialogue, it assesses response helpfulness using the dataset's preferred responses as the baseline. Prior work suggests that language models may outperform existing automated metrics as evaluators \citep{Chen2023ExploringTU}; nonetheless, we perform a human-subjects study, described in Sec.~\ref{sec:human-judgments}, to validate the application of GPT-4 for this purpose. Our findings demonstrate strong alignment between GPT-4 and human assessments, with human agreement with GPT-4 rivaling or surpassing the level of consistency observed among different human annotators.

\textbf{Methods.} Alongside DPO, we benchmark several established methods for training language models to reflect human preferences. For the summarization task, we consider zero-shot prompting with \textbf{GPT-J} \citep{gpt-j}, and for dialogue, 2-shot prompting with \textbf{Pythia-2.8B} \citep{biderman2023pythia}. We further examine the \textbf{SFT} approach, and the \textbf{Preferred-FT} variant, which involves supervised fine-tuning on the selected response $y_w$—obtained from either the SFT model (for sentiment and summarization tasks) or from a general-purpose language model (for dialogue). Another technique, \textbf{Unlikelihood}~\citep{welleck2019neural}, promotes the policy's likelihood of $y_w$ while actively suppressing the likelihood of $y_l$; we include an adjustable coefficient $\alpha \in [0,1]$ for the unlikelihood component. Our comparisons extend to \textbf{PPO} \citep{schulman2017proximal}, utilizing a reward function inferred from preference data, as well as to \textbf{PPO-GT}, which acts as an oracle by employing the true reward function in the controlled sentiment setting. For sentiment-focused experiments, two PPO-GT variants are assessed: a standard off-the-shelf implementation \cite{leandro_von_werra_2023_7790115} and an enhanced version incorporating reward normalization and additional hyperparameter tuning (these improvements are also applied when training PPO with learned rewards). Lastly, we include the \textbf{Best of $N$} approach as a baseline, in which $N$ outputs are sampled from the SFT model (or Preferred-FT for dialogue) and the candidate with the highest score—according to a reward model trained on preference data—is selected. While this strategy effectively isolates the impact of the reward model from PPO-based optimization, its practicality is limited by the computational cost of generating $N$ samples per input at test time, even for moderate values of $N$.

\subsection{How well can DPO optimize the RLHF objective?}

In standard RLHF approaches, the KL-constrained reward maximization objective is designed to maximize reward while preventing the policy from straying significantly from the reference policy. Consequently, a meaningful comparison between algorithms must consider both the obtained reward and the KL divergence; attaining marginally greater reward at the cost of substantially increased KL is often suboptimal. In Figure~\ref{fig:frontier-tldr-main}, we illustrate the trade-off curve between reward and KL for different algorithms on the sentiment task. For each method, we conduct multiple independent training runs, each using a distinct parameter to modulate policy conservativeness (for PPO, target KL values are $\in\{3,6,9,12\}$; for DPO, $\beta$ spans $\{0.05,0.1,1,5\}$; unlikelihood employs $\alpha\in\{0.05,0.1,0.5,1\}$; and for preferred-FT, we utilize different random seeds). This comprehensive sweep covers 22 total experiments. At intervals of 100 training steps until convergence, we evaluate the learned policies on a set of held-out prompts, computing both the mean reward under the true reward model and the mean sequence-level KL\footnote{This is calculated as the sum of KL-divergences at each timestep.} with respect to the reference policy, denoted $\text{KL}\left(\pi\mid \mid \pi_\text{ref}\right)$. Our findings indicate that DPO yields a Pareto frontier that is markedly superior, achieving the highest rewards while maintaining low KL. This outcome is noteworthy for several reasons: Firstly, even though DPO and PPO are designed to optimize the identical objective, DPO demonstrates significantly greater efficiency—its reward/KL tradeoff entirely outperforms PPO. Secondly, DPO surpasses PPO’s frontier \textit{even in scenarios where PPO is provided with access to exact rewards} (i.e., PPO-GT).

\subsection{Can DPO scale to real preference datasets?}
\label{sec:dpo-real-datasets}
We proceed to assess the efficacy of DPO for fine-tuning on tasks involving summarization and single-turn dialogue. In summarization, standard automatic metrics like ROUGE are known to align poorly with human judgment~\citep{stiennon2022learning}; previous research has demonstrated improvements in summary quality when LMs are fine-tuned with PPO using human preference data. To benchmark various approaches, we generate outputs using the test set from the TL;DR summarization dataset, measuring the mean win rate against reference outputs from the test split. Output samples for each method are generated at a range of temperatures from 0.0 to 1.0, and the corresponding win rates are illustrated in Figure~\ref{fig:frontier-tldr-main} (right). All methods—including DPO, PPO, and Preferred-FT—are applied to the same GPT-J SFT model\footnote{\url{https://huggingface.co/CarperAI/openai_summarize_tldr_sft}} for fine-tuning. Experimental results indicate that DPO achieves a win rate near 61\% at a temperature of 0.0, surpassing PPO's win rate of approximately 57\% at its best-performing temperature, also 0.0. DPO further establishes a higher peak win rate relative to the best of $N$ baseline. It is important to mention that DPO's $\beta$ hyperparameter was not systematically optimized, suggesting that the current performance may not fully reflect its potential. Additionally, DPO demonstrates substantially greater stability to variations in sampling temperature compared to PPO, whose win rate can decline to the level of the original GPT-J model at elevated temperatures. The Preferred-FT approach does not yield notable improvements over the SFT baseline. Furthermore, Section~\ref{sec:human-judgments} presents a direct comparison between DPO and PPO using human assessments, where DPO samples at a temperature of 0.25 were favored 58\% of the time in comparison to PPO outputs at the same temperature.

For single-turn dialogue, we assess various approaches using a subset of the test portion from the Anthropic HH dataset \citep{bai2022training} that contains one round of human-assistant interaction. In our GPT-4 assessments, the preferred completions on the test set serve as the ground truth for calculating the win rates of different techniques. As there is no widely adopted SFT baseline for this setting, our experiments commence with a pre-trained Pythia-2.8B model; we apply Preferred-FT to create a reference model that is finetuned on the selected completions to ensure the generated outputs remain in-distribution, followed by training with DPO. Additional baselines include the best outcome from 128 Preferred-FT completions (as the performance of the Best of $N$ method saturates at $N=128$ for this particular task; details in Appendix Figure~\ref{fig:best-of-n}), as well as a 2-shot prompt adaptation of the Pythia-2.8B base model. Our findings indicate that DPO matches or exceeds the performance of other methods at their optimal sampling temperatures. Furthermore, we test an RLHF model trained via PPO on the Anthropic HH dataset\footnote{\url{https://huggingface.co/reciprocate/ppo_hh_pythia-6B}} sourced from a prominent implementation\footnote{\url{https://github.com/CarperAI/trlx/tree/main/examples/hh}}, but cannot identify any prompting strategy or temperature yielding results superior to the original Pythia-2.8B model. Drawing on insights from TL;DR, and the observation that both PPO and the Best of 128 method maximize the same reward objective, we regard Best of 128 as an approximate upper bound for PPO-level performance. Notably, among the evaluated strategies, DPO emerges as the only computationally efficient approach that outperforms the preferred completions in the Anthropic HH test set and achieves results that are comparable to, or surpass, those of the resource-intensive Best of 128 method. Finally, as depicted in Figure~\ref{fig:dialogue-main}, DPO reaches its optimal performance in relatively few training iterations.

\subsection{Generalization to a new input distribution}

To more rigorously assess how PPO and DPO perform when exposed to distribution shifts, we evaluated the policies obtained from our Reddit TL;DR summarization experiments on a distinct dataset: news articles drawn from the test portion of the CNN/DailyMail dataset \citep{nallapati-etal-2016-abstractive}. These evaluations used the optimal sampling temperatures determined from the TL;DR setting (0 and 0.25). The corresponding outcomes are detailed in Table~\ref{tab:ood}. To assess summary quality, we computed the win rate of GPT-4 against the reference summaries for both datasets, employing the same GPT-4 (C) prompt as used for Reddit TL;DR—modifying only the phrase ``forum post'' to ``news article''. Within this out-of-distribution context, DPO maintains a sizable performance advantage over PPO. This result constitutes preliminary evidence that DPO policies retain strong generalization properties, matching or exceeding those of PPO, despite DPO's not relying on the extra unlabeled Reddit TL;DR prompts leveraged by PPO.

\subsection{Validating GPT-4 judgments with human judgments}
\label{sec:human-judgments}
We conducted a human evaluation study to ascertain the alignment between GPT-4's assessments and those of human annotators. This involved examining outputs from the TL;DR summarization experiment under two different GPT-4 prompting strategies. The \textbf{GPT-4 (S)} (simple) prompt directly asks which summary better encapsulates the essential information from the post. In contrast, the \textbf{GPT-4 (C)} (concise) prompt explicitly instructs GPT-4 to also consider conciseness; this distinction was motivated by observations that GPT-4, under the (S) prompt, tends to favor more verbose and repetitive summaries compared to human preferences. Complete prompt formulations are provided in Appendix~\ref{app:prompts}. Three comparative evaluations were performed, corresponding to the top-performing method (DPO, temp. 0.25), the lowest-performing method (PPO, temp. 1.0), and an intermediate method (SFT, temp. 0.25), to ensure a spectrum of output qualities; each method was contrasted against PPO with greedy decoding (the temperature yielding its highest performance). Our analysis shows that for both prompts, GPT-4’s preferences are as consistent with those of human raters as inter-rater reliability among humans themselves, supporting GPT-4's utility as an effective proxy for human evaluations. Due to resource limitations, multiple human judgments were collected only for the DPO and PPO-1 comparisons. In general, the \textbf{GPT-4 (C)} prompt produces win rates that are more reflective of human assessments, and for this reason, we adopt it for principal reporting in Section~\ref{sec:dpo-real-datasets}. Further methodological details about the human study, such as the web interface used for rating and the roster of human participants, can be found in Appendix~\ref{app:human-study}.

\section{Discussion}
Preference-based learning offers a robust and adaptable approach for developing proficient and value-aligned language models. In this work, we presented DPO, an uncomplicated training method that enables language models to learn from preference data without employing reinforcement learning. Instead of reframing preference learning within the conventional reinforcement learning paradigm to leverage standard RL algorithms, DPO establishes a direct correspondence between language model policies and reward functions. This allows the model to be trained to align with human preferences using a straightforward cross-entropy loss function, thereby circumventing reinforcement learning and maintaining broad applicability. DPO achieves competitive or superior performance to established RLHF techniques, such as those utilizing PPO, and requires minimal hyperparameter adjustment. As a result, DPO lowers the threshold for preference-based training of language models.

\textbf{Limitations \& Future Work.} Our findings invite several avenues for further investigation. One key question is how DPO-trained policies handle out-of-distribution scenarios compared to approaches that rely on explicit reward modeling. Preliminary results imply that the generalization capabilities of DPO are on par with PPO-trained counterparts, but thorough empirical validation remains necessary. For instance, future work could examine whether self-labeling techniques using DPO policies can leverage unlabeled data as effectively. Another area of interest involves understanding the manifestation of reward over-optimization within DPO, and whether the marginal decline in performance observed in Figure~\ref{fig:dialogue-main}-right can be attributed to this phenomenon. Furthermore, while our current evaluation extends to models with up to 6B parameters, scaling DPO to contemporary models that are orders of magnitude larger represents a promising direction. In terms of assessment methodologies, we observe that the win rates judged by GPT-4 are sensitive to prompt phrasing; investigating optimal prompting strategies for robust automated evaluation remains an open question. Lastly, the scope of DPO extends beyond language models trained from human preferences, offering opportunities for its application in the training of generative models across different data modalities.